\documentclass[a4paper,12pt,oneside]{maki-notes}

\renewcommand{\LectureNotesTitle}{TikZ 模板页示例}
\renewcommand{\AuthorInfo}{Maki Notes Template Pack}
\renewcommand{\CoverTitle}{TikZ 模板页示例}
\renewcommand{\CoverAuthorBlock}{%
	Maki Notes 图形模板集
	\\[0.8em]
	\Large
	Transformer, CNN, ResNet, 球面, 流形, 傅里叶变换, 交换图, 流程图, 时间轴
	\\
}
\renewcommand{\CoverFooterBlock}{%
	参考 tikz.net 的结构思路改写为讲义模板页
	\\[0.5em]
	适合直接复制到课程讲义、论文附录与技术报告
	\\[2.5em]
	二〇二六年四月
}

\begin{document}

\frontmatter
\MakeTitlePage

\chapter*{使用说明}
这一份文档不是测试页, 而是现成模板页集合. 每个小节都对应一个可以直接复制进正式文档的图版: 你可以改标签文字、层数、颜色或注释框, 但尽量保留节点命名和样式名称, 这样后续维护成本最低.

本文件的图形结构主要参考了 tikz.net 上的若干页面, 但已经改写成适合本模板的语义样式与讲义版面:
\begin{itemize}
\item \texttt{neural\_networks} 提供了网络层与张量块的组织思路
\item \texttt{skip-connection} 提供了残差支路的结构思路
\item \texttt{manifold} 提供了流形与潜空间映射的版面思路
\item \texttt{parametric-sphere} 提供了球面、经纬线与切平面的结构思路
\item \texttt{fourier-transform} 提供了时域/频域并排展示的结构思路
\item \texttt{plane-to-torus} 与 \texttt{torus} 提供了拓扑识别与曲面表达的思路
\item \texttt{flowchart} 与 \texttt{timeline} 提供了信息流程和时间组织的版面思路
\item \texttt{venn-diagrams} 与 \texttt{graph-isomorphism} 提供了概率/集合图与节点边图的组织思路
\item 交换图部分则基于当前模板已有的 \texttt{diag} 语义样式继续扩展
\end{itemize}

\tableofcontents

\mainmatter

\chapter{深度学习模板页}

\section{Transformer 编码器模板}

这一页适合放在“模型结构”或“序列建模”一章. 版式重点不是复刻论文中的每一个算子, 而是把 token、注意力、前馈层、残差与归一化的逻辑顺序表达清楚.

\begin{center}
\begin{tikzpicture}[x=0.79cm,y=0.88cm]
\foreach \i/\txt in {1/$x_1$,2/$x_2$,3/$x_3$,4/$x_4$,5/$x_5$}{
  \node[nn token] (tok\i) at (\i*1.4,0) {\txt};
  \draw[nn attention] (tok\i.north) -- ++(0,0.7);
}

\node[nn op accent, minimum width=8.2cm, minimum height=1cm] (embed) at (4.2,1.5) {Embedding $+$ Positional Encoding};
\foreach \i in {1,...,5}{
  \draw[nn connect strong] (tok\i.north) -- (embed.south -| tok\i.north);
}

\node[nn op warm, minimum width=8.2cm, minimum height=1.15cm] (attn) at (4.2,3.2) {Multi-Head Self-Attention};
\node[nn sum] (add1) at (4.2,4.5) {$+$};
\node[nn op, minimum width=8.2cm, minimum height=0.95cm] (norm1) at (4.2,5.8) {LayerNorm};
\node[nn op warm, minimum width=8.2cm, minimum height=1.15cm] (ffn) at (4.2,7.5) {Feed-Forward Network};
\node[nn sum] (add2) at (4.2,8.8) {$+$};
\node[nn op, minimum width=8.2cm, minimum height=0.95cm] (norm2) at (4.2,10.1) {LayerNorm};

\draw[nn connect strong] (embed.north) -- (attn.south);
\draw[nn connect strong] (attn.north) -- (add1.south);
\draw[nn connect strong] (add1.north) -- (norm1.south);
\draw[nn connect strong] (norm1.north) -- (ffn.south);
\draw[nn connect strong] (ffn.north) -- (add2.south);
\draw[nn connect strong] (add2.north) -- (norm2.south);

\draw[nn skip] ($(embed.east)+(0.15,0)$) .. controls +(1.2,1.1) and +(1.2,-1.1) .. ($(add1.east)+(0.15,0)$);
\draw[nn skip] ($(norm1.west)+(-0.15,0)$) .. controls +(-1.2,1.1) and +(-1.2,-1.1) .. ($(add2.west)+(-0.15,0)$);

\foreach \i/\txt in {1/$h_1$,2/$h_2$,3/$h_3$,4/$h_4$,5/$h_5$}{
  \node[nn token] (out\i) at (\i*1.4,11.6) {\txt};
  \draw[nn connect strong] (norm2.north -| out\i.south) -- (out\i.south);
}

\node[maki annotation] at (8.9,3.2) {attention\\mixes token\\interactions};
\node[maki annotation] at (-0.6,6.7) {residual\\paths};
\node[maki annotation] at (8.9,7.5) {position-wise\\nonlinearity};
\end{tikzpicture}
\captionof{figure}{Transformer 编码器块模板}
\end{center}

\section{CNN 特征提取模板}

这一页适合在卷积神经网络、视觉表征或多尺度特征章节里使用. 结构参考 tikz.net 的 \texttt{neural\_networks} 中 CNN 表达方式, 但这里改成了更适合课程讲义的“张量尺寸 + 算子流程”版.

\begin{center}
\begin{tikzpicture}[x=0.9cm,y=1cm]
\node[nn tensor, minimum width=1.8cm, minimum height=2.6cm, label={[maki label]below:$224\times224\times3$}] (img) at (0,0) {image};

\foreach \x/\w/\h/\lab/\dim in {
2.3/0.45/1.8/{Conv\\$3\times3$}/{112\times112\times64},
3.0/0.55/2.2/{Conv\\$3\times3$}/{112\times112\times64},
4.4/0.7/2.7/{Pool}/{56\times56\times64},
6.0/0.8/3.0/{Conv\\Block}/{56\times56\times128},
7.5/0.95/3.35/{Conv\\Block}/{28\times28\times256},
9.1/1.05/3.6/{Global\\Pool}/{1\times1\times256}
}{
  \draw[nn feature map, fill=NNTensorColor!10] (\x-\w/2,-\h/2) rectangle ++(\w,\h);
  \draw[nn feature map, fill=NNTensorColor!18] (\x-\w/2+0.16,-\h/2+0.16) rectangle ++(\w,\h);
  \draw[nn feature map, fill=NNTensorColor!26] (\x-\w/2+0.32,-\h/2+0.32) rectangle ++(\w,\h);
  \node[maki label] at (\x,\h/2+0.95) {\lab};
  \node[maki label] at (\x,-\h/2-0.7) {$\dim$};
}

\node[nn output, minimum size=11mm] (cls) at (11.3,0.2) {$\hat y$};

\draw[nn connect strong] (img.east) -- (1.75,0);
\draw[nn connect strong] (1.75,0) -- (2.25,0);
\foreach \a/\b in {2.65/3.35,3.35/5.0,5.0/6.8,6.8/8.3,8.3/10.6}{
  \draw[nn connect strong] (\a,0) -- (\b,0);
}
\draw[nn connect strong] (10.95,0.2) -- (cls.west);

\draw[nn attention] (3.0,2.0) -- +(0.95,0.75) node[maki annotation, above right] {shared kernels};
\draw[nn attention] (7.5,2.55) -- +(1.05,0.75) node[maki annotation, above right] {deeper semantics};
\node[maki annotation] at (11.3,1.8) {classifier};
\end{tikzpicture}
\captionof{figure}{CNN 特征提取与分类头模板}
\end{center}

\section{ResNet 残差块模板}

这一页可以直接用于介绍残差学习的动机. 它参考了 tikz.net 的 \texttt{skip-connection} 结构, 但这里改成了更适合板书式讲义的块图写法.

\begin{center}
\begin{tikzpicture}[x=0.86cm,y=1cm]
\node[nn tensor, minimum width=1.5cm, minimum height=0.9cm] (x) at (0,0) {$\bm x$};
\node[nn op warm] (conv1) at (2.3,0) {Conv\\BN\\ReLU};
\node[nn op warm] (conv2) at (5.0,0) {Conv\\BN};
\node[nn sum] (add) at (7.5,0) {$+$};
\node[nn op accent] (relu) at (9.2,0) {ReLU};
\node[nn tensor, minimum width=1.5cm, minimum height=0.9cm] (y) at (11.1,0) {$\bm y$};

\draw[nn connect strong] (x.east) -- (conv1.west);
\draw[nn connect strong] (conv1.east) -- (conv2.west);
\draw[nn connect strong] (conv2.east) -- (add.west);
\draw[nn connect strong] (add.east) -- (relu.west);
\draw[nn connect strong] (relu.east) -- (y.west);

\draw[nn skip] (x.north) .. controls +(0,1.8) and +(0,1.8) .. (add.north);
\node[maki annotation] at (3.7,2.0) {residual branch\\$\mathcal F(\bm x)$};
\node[maki annotation] at (6.2,-1.9) {identity map};

\draw[nn residual bracket] (1.4,-1.2) -- (1.4,-1.55) -- (6.0,-1.55) -- (6.0,-1.2);
\node[maki label] at (3.7,-2.05) {residual block};

\node[nn op, minimum width=2.1cm] (proj) at (3.0,-3.55) {Proj.\ skip};
\draw[nn connect strong] ($(x.south)+(0,-0.1)$) -- ++(0,-0.65) -| (proj.west);
\draw[nn skip] (proj.east) -| ($(add.south)+(0,-0.1)$);
\node[maki annotation] at (8.8,-3.55) {optional $1\times1$ conv when channel\\or stride mismatch occurs};
\end{tikzpicture}
\captionof{figure}{ResNet 残差块模板}
\end{center}

\chapter{数学与拓扑模板页}

\section{拓扑交换图模板}

这一页适合范畴论、代数拓扑或微分几何中的自然交换图. 它不是 `tikz-cd` 的替代品, 而是为了和当前样式系统保持统一色彩与盒注释风格.

\begin{center}
\begin{tikzpicture}[x=2.7cm,y=1.7cm]
\node[diag object] (A) at (0,2) {$\pi_1(X,x_0)$};
\node[diag object] (B) at (2.4,2) {$H_1(X;\mathbb Z)$};
\node[diag object] (C) at (0,0) {$\pi_1(Y,f(x_0))$};
\node[diag object] (D) at (2.4,0) {$H_1(Y;\mathbb Z)$};
\node[diag object] (E) at (4.9,1) {$\operatorname{Hom}(G,\mathbb Z)$};

\draw[diag morphism] (A) -- node[above] {$\operatorname{ab}_X$} (B);
\draw[diag morphism] (C) -- node[below] {$\operatorname{ab}_Y$} (D);
\draw[diag morphism emph] (A) -- node[left] {$f_{\ast}$} (C);
\draw[diag morphism emph] (B) -- node[right] {$f_{\ast}$} (D);
\draw[diag morphism] (B) -- node[above right] {$\Phi_X$} (E);
\draw[diag morphism] (D) -- node[below right] {$\Phi_Y$} (E);
\draw[diag morphism dashed] (A) -- node[above left] {$\chi$} (E);
\draw[diag morphism dashed] (C) -- node[below left] {$\chi\circ f_{\ast}$} (E);

\node[maki annotation] at (1.2,1.0) {left square commutes};
\node[maki annotation] at (3.8,2.0) {universal\\property};
\end{tikzpicture}
\captionof{figure}{代数拓扑/范畴论中的交换图模板}
\end{center}

\clearpage

\section{交换图与 pullback 模板}

这一页适合需要同时表达 pullback、自然变换或“背景范畴对象”的场景. 它沿用当前模板已有的 `diag` 系列样式, 但补上了强调对象、淡化对象、弯箭头和二维单元标记, 更适合讲交换图的“结构角色”而不是只画箭头网格.

\begin{center}
\begin{tikzpicture}[x=2.35cm,y=1.75cm]
\node[diag object accent] (A) at (0,2) {$X\times_Z Y$};
\node[diag object] (B) at (2.2,2) {$Y$};
\node[diag object] (C) at (0,0) {$X$};
\node[diag object accent] (D) at (2.2,0) {$Z$};
\node[diag object muted] (E) at (4.5,1) {$\mathcal C$};

\draw[diag morphism emph] (A) -- node[above] {$p_2$} (B);
\draw[diag morphism emph] (A) -- node[left] {$p_1$} (C);
\draw[diag morphism] (B) -- node[right] {$g$} (D);
\draw[diag morphism] (C) -- node[below] {$f$} (D);
\draw[diag morphism dotted] (A) to[bend left=10] node[above left] {$u$} (E);
\draw[diag morphism natural] (B) -- node[above right] {$G$} (E);
\draw[diag morphism natural] (C) -- node[below right] {$F$} (E);
\draw[diag morphism curve] (A) to[bend right=18] node[below left] {$\eta$} (D);
\draw[diag corner] ($(A.south east)+(0.16,-0.02)$) -- ++(0.22,0) -- ++(0,-0.22);
\node[diag 2cell] at (1.1,1.0) {pb};

\node[maki annotation] at (2.2,2.8) {pullback square};
\node[maki annotation] at (4.5,-0.12) {ambient / derived\\context object};
\end{tikzpicture}
\captionof{figure}{pullback、自然变换与强调对象的交换图模板}
\end{center}

\clearpage

\section{三维球面与切平面模板}

这一页适合微分几何、向量分析或流形导论中介绍球面、经纬线和切平面. 图形组织参考了 tikz.net 的 \texttt{parametric-sphere} 页面, 但收敛成更适合讲义版面的轻量模板: 保留可见/不可见轮廓、经纬线与切平面的层次, 不追求过重的三维渲染.

\begin{center}
\begin{tikzpicture}[x=1cm,y=1cm]
\shade[sphere shell] (0,0) circle (2.0);
\draw[sphere hidden] (-2.0,0) arc[start angle=180,end angle=360,x radius=2.0,y radius=0.62];
\draw[sphere equator] (2.0,0) arc[start angle=0,end angle=180,x radius=2.0,y radius=0.62];
\draw[sphere hidden] (0,2.0) arc[start angle=90,end angle=270,x radius=0.8,y radius=2.0];
\draw[sphere longitude] (0,-2.0) arc[start angle=-90,end angle=90,x radius=0.8,y radius=2.0];
\draw[sphere latitude] (-1.55,1.05) arc[start angle=180,end angle=360,x radius=1.55,y radius=0.36];
\draw[sphere latitude] (1.35,-0.95) arc[start angle=0,end angle=180,x radius=1.35,y radius=0.28];
\filldraw[sphere tangent plane] (1.2,1.05) -- (2.5,1.45) -- (2.1,2.25) -- (0.85,1.85) -- cycle;
\node[sphere point, label={[maki label]above right:$p$}] (sp) at (1.55,1.35) {};
\draw[sphere vector] (0,0) -- (sp);
\draw[sphere axis] (0,0) -- (2.8,0) node[right] {$x$};
\draw[sphere axis] (0,0) -- (0,2.7) node[above] {$z$};
\node[maki annotation] at (0,-2.65) {sphere shell\\latitude / longitude\\tangent plane};
\node[maki annotation] at (2.7,2.35) {$T_pS^2$};
\end{tikzpicture}
\captionof{figure}{球面、经纬线与切平面模板}
\end{center}

\clearpage

\section{局部流形图与坐标图模板}

这一页更偏微分几何或“局部欧氏化”表达. 它与后面的“流形与潜空间”模板不同, 重点不是编码器/潜变量, 而是 patch、测地线、局部 chart 和点的坐标映射关系.

\begin{center}
\begin{tikzpicture}[x=1cm,y=1cm]
\draw[manifold patch] (0.2,1.05) .. controls (1.05,2.28) and (2.3,2.6) .. (3.4,2.08)
  .. controls (4.3,1.66) and (4.75,0.88) .. (4.42,0.12)
  .. controls (3.68,-0.45) and (2.05,-0.2) .. (1.0,0.28)
  .. controls (0.48,0.52) and (0.12,0.84) .. cycle;
\draw[manifold geodesic] (0.82,0.92) .. controls (1.45,1.78) and (2.45,2.02) .. (3.42,1.76)
  .. controls (4.02,1.62) and (4.35,1.22) .. (4.18,0.72);
\filldraw[manifold tangent] (2.12,1.08) -- (3.12,1.38) -- (2.82,2.14) -- (1.88,1.84) -- cycle;
\draw[manifold fiber] (2.48,1.62) -- ++(0,1.35);
\node[manifold point, label={[maki label]above:$q$}] (mq) at (2.48,1.62) {};

\draw[manifold chart] (5.32,1.02) rectangle (6.88,2.56);
\draw[topo mesh] (5.84,1.02) -- (5.84,2.56);
\draw[topo mesh] (6.36,1.02) -- (6.36,2.56);
\draw[topo mesh] (5.32,1.53) -- (6.88,1.53);
\draw[topo mesh] (5.32,2.04) -- (6.88,2.04);
\draw[manifold map accent] (mq) to[out=16,in=170] node[midway,above] {$\varphi$} (5.92,1.82);
\node[maki label] at (5.92,1.82) {$\varphi(q)$};

\node[maki annotation] at (2.2,-0.72) {patch / geodesic / tangent plane};
\node[maki annotation] at (6.12,3.02) {local chart};
\end{tikzpicture}
\captionof{figure}{局部流形图、测地线与坐标图模板}
\end{center}

\clearpage

\section{流形与潜空间示意模板}

这一页参考 tikz.net 的 \texttt{manifold} 页面思路, 适合写“高维数据流形假设”“表示学习”“自编码器几何解释”之类内容.

\begin{center}
\begin{tikzpicture}[x=1cm,y=1cm]
\draw[manifold axis] (0,0) -- (0,2.8);
\draw[manifold axis] (0,0) -- (2.8,0.75);
\draw[manifold axis] (0,0) -- (3.5,0) node[midway,below=4pt] {ambient space $\mathcal X$};

\fill[manifold cloud]
  (0.35,0.85) to[out=10,in=155] (2.1,0.45)
  to[out=25,in=210] (3.25,1.4)
  to[out=125,in=15] (2.35,2.15)
  to[out=190,in=40] (0.9,2.0) to[out=220,in=135] cycle;

\shade[left color=TopologyBg,right color=TopologyColor!28,draw=none]
  (0.55,0.95) to[out=5,in=150] (2.0,0.55)
  to[out=22,in=210] (3.0,1.25)
  to[out=130,in=10] (2.2,1.95)
  to[out=195,in=32] (1.05,1.8) to[out=220,in=135] cycle;

\draw[manifold axis] (6.0,0.8) -- (6.0,3.0);
\draw[manifold axis] (6.0,0.8) -- (8.25,0.8) node[midway,below=4pt] {latent space $\mathcal F$};
\draw[manifold latent] (6.65,1.25) rectangle (8.0,2.6);
\draw[topo mesh] (7.1,1.25) -- (7.1,2.6);
\draw[topo mesh] (7.55,1.25) -- (7.55,2.6);
\draw[topo mesh] (6.65,1.7) -- (8.0,1.7);
\draw[topo mesh] (6.65,2.15) -- (8.0,2.15);

\draw[manifold forward] (2.1,1.45) to[out=25,in=165] node[midway,above] {$f$} (6.75,2.0);
\draw[manifold backward] (7.85,1.8) to[out=205,in=18] node[midway,below] {$g$} (2.65,1.1);

\node[geom highlighted point, label={[maki label]above:$x$}] at (1.55,1.45) {};
\node[geom highlighted point, label={[maki label]above:$x'$}] at (2.55,1.55) {};
\node[geom highlighted point, label={[maki label]right:$z$}] at (7.2,2.0) {};
\draw[topo geodesic] (1.55,1.45) .. controls (1.8,1.9) and (2.25,1.9) .. (2.55,1.55);
\node[maki annotation] at (2.2,2.45) {data manifold};
\node[maki annotation] at (7.35,3.0) {chart / code};
\end{tikzpicture}
\captionof{figure}{流形假设与潜空间映射模板}
\end{center}

\clearpage

\section{傅里叶变换模板}

这一页适合分析、信号处理或偏微分方程讲义里快速说明“时域到频域”的关系. 它参考了 tikz.net 的 \texttt{fourier-transform} 页面, 但把常用元素拆成了时域信号、频域谱线、窗函数标注和变换箭头, 方便直接替换成矩形窗、高斯或离散谱的版本.

\begin{center}
\begin{tikzpicture}[x=1cm,y=1cm]
\begin{scope}
  \draw[fourier axis] (-3.4,0) -- (3.45,0) node[right] {$t$};
  \draw[fourier axis] (0,-0.45) -- (0,1.95) node[above] {$f(t)$};
  \draw[fourier signal] (-3.0,0) -- (-1.1,0) -- (-1.1,1.18) -- (1.1,1.18) -- (1.1,0) -- (3.05,0);
  \draw[fourier guide] (-1.1,0) -- (-1.1,1.18);
  \draw[fourier guide] (1.1,0) -- (1.1,1.18);
  \node[fourier marker, label={[maki label]below:$-T$}] at (-1.1,0) {};
  \node[fourier marker, label={[maki label]below:$T$}] at (1.1,0) {};
  \node[maki annotation] at (0.0,1.72) {time domain};
\end{scope}

\draw[fourier arrow] (4.25,0.95) -- (6.0,0.95) node[midway,above] {$\mathcal F$};

\begin{scope}[xshift=10cm]
  \draw[fourier axis] (-3.3,0) -- (3.4,0) node[right] {$\omega$};
  \draw[fourier axis] (0,-0.55) -- (0,2.0) node[above] {$\hat f(\omega)$};
  \draw[fourier spectrum, samples=180, smooth, variable=\x, domain=-3.0:3.0]
    plot (\x,{1.3*(abs(\x)<0.06 ? 1 : sin(deg(2.2*\x))/(2.2*\x))});
  \foreach \x/\lbl in {-2.92/$-\frac{2\pi}{T}$,-1.46/$-\frac{\pi}{T}$,1.46/$\frac{\pi}{T}$,2.92/$\frac{2\pi}{T}$}{
    \draw[fourier guide] (\x,0) -- (\x,{1.3*(abs(\x)<0.06 ? 1 : sin(deg(2.2*\x))/(2.2*\x))});
    \node[fourier marker] at (\x,0) {};
    \node[maki label, below] at (\x,-0.04) {\lbl};
  }
  \node[maki annotation] at (0.15,1.8) {frequency domain};
  \node[fourier window, anchor=west] at (1.18,1.24) {$\hat f(\omega)\propto \dfrac{\sin(T\omega)}{\omega}$};
\end{scope}
\end{tikzpicture}
\captionof{figure}{矩形窗的时域/频域变换模板}
\end{center}

\clearpage

\section{环面与基本域模板}

这一页适合在拓扑、微分几何或群作用章节中解释“边识别”和“商空间”. 这里把基本域与环面的概念图并排放置, 比单独画其中一个更适合教学.

\begin{center}
\begin{tikzpicture}[x=1cm,y=1cm]
\begin{scope}
  \draw[topo patch] (0,0) rectangle (4.4,4.4);
  \foreach \k in {0.73,1.46,2.19,2.92,3.65}{
    \draw[topo mesh] (\k,0) -- (\k,4.4);
    \draw[topo mesh] (0,\k) -- (4.4,\k);
  }
  \draw[topo identify] (0,0) -- (4.4,0);
  \draw[topo identify] (0,4.4) -- (4.4,4.4);
  \draw[topo identify reversed] (0,0) -- (0,4.4);
  \draw[topo identify reversed] (4.4,0) -- (4.4,4.4);
  \draw[topo loop] (0.9,1.0) .. controls (1.6,2.8) and (2.8,2.7) .. (3.7,1.3);
  \draw[topo geodesic] (0.7,3.4) .. controls (1.7,3.9) and (3.1,3.0) .. (3.8,3.55);
  \node[maki annotation] at (2.2,2.2) {fundamental\\domain};
  \node[maki label] at (2.2,-0.45) {$a$};
  \node[maki label] at (2.2,4.85) {$a$};
  \node[maki label] at (-0.35,2.2) {$b^{-1}$};
  \node[maki label] at (4.75,2.2) {$b^{-1}$};
\end{scope}

\begin{scope}[xshift=7.6cm]
  \draw[topo boundary, fill=TopologyBg] (0,0) ellipse (2.25 and 1.38);
  \draw[topo boundary, fill=white] (0,0) ellipse (0.84 and 0.46);
  \draw[topo loop] (0,0.72) ellipse (1.8 and 0.42);
  \draw[topo geodesic] (-1.35,-0.08) .. controls (-0.45,1.18) and (1.0,1.08) .. (1.45,0.02);
  \draw[topo cut] (0,-1.38) -- (0,1.38);
  \draw[maki guide] (-2.35,0) -- (2.35,0);
  \draw[maki guide] (0,-1.55) -- (0,1.55);
  \node[maki annotation] at (0,-2.0) {identified\\torus};
\end{scope}
\begin{scope}[xshift=7.6cm]
  \draw[manifold forward] (-3.0,0.2) to[out=15,in=170] node[midway,above] {quotient} (-2.1,0.25);
\end{scope}
\end{tikzpicture}
\captionof{figure}{基本域到环面的拓扑模板}
\end{center}

\chapter{通用讲义图形模板页}

\section{流程图模板}

这一页适合放在“算法流程”“证明路线”“实验流水线”或“数据处理步骤”小节里. 重点是把开始节点、处理步骤、判断节点、旁路反馈和补充说明分成独立样式, 这样后续改流程长度时不需要重写所有参数.

\begin{center}
\begin{tikzpicture}[x=1cm,y=1cm]
\node[flow terminal] (start) at (0,2.2) {开始};
\node[flow process] (collect) at (2.8,2.2) {收集输入\\与设定参数};
\node[flow process] (build) at (5.8,2.2) {构造模型\\或中间对象};
\node[flow decision] (check) at (8.8,2.2) {满足\\约束?};
\node[flow process] (polish) at (11.9,2.2) {整理输出\\与可视化};
\node[flow terminal] (finish) at (14.8,2.2) {结束};

\node[flow io] (revise) at (8.8,-0.4) {调整参数\\补充条件};
\node[flow io] (archive) at (11.9,-0.4) {保存日志\\导出结果};
\node[flow note, anchor=west] at (12.0,-0.4) {把流程说明、例外分支和\\导出动作拆开写, 适合直接用于\\课程讲义或实验报告.};

\draw[flow arrow] (start) -- (collect);
\draw[flow arrow] (collect) -- (build);
\draw[flow arrow] (build) -- (check);
\draw[flow arrow emph] (check) -- node[above, maki label] {是} (polish);
\draw[flow arrow] (polish) -- (finish);
\draw[flow arrow] (polish) -- (archive);
\draw[flow arrow feedback] (check) -- node[left, maki label] {否} (revise);
\draw[flow arrow] (revise.east) -| (build.south);

\node[maki annotation] at (2.8,4.0) {准备阶段};
\node[maki annotation] at (5.8,4.0) {主流程};
\node[maki annotation] at (11.9,4.0) {输出阶段};
\end{tikzpicture}
\captionof{figure}{讲义/算法说明中的流程图模板}
\end{center}

\clearpage

\section{时间轴模板}

这一页适合课程安排、项目里程碑、论文结构导读或者历史脉络概览. 如果你想表达的是“阶段 + 节点”而不是普通折线图, 这类时间轴比表格更容易在讲义中快速建立节奏感.

\begin{center}
\begin{tikzpicture}[x=0.92cm,y=0.92cm]
\draw[timeline axis] (0.2,0) -- (14.5,0);
\foreach \x/\label in {1.0/第 1 周,3.2/第 3 周,5.5/第 5 周,8.1/第 8 周,10.9/第 11 周,13.4/第 14 周}{
  \draw[timeline tick] (\x,-0.18) -- (\x,0.18);
  \node[maki label, below] at (\x,-0.2) {\label};
}

\draw[timeline span] (1.2,0.85) -- (5.2,0.85);
\draw[timeline span, draw=LogicColor!74!black] (5.8,0.85) -- (10.6,0.85);
\draw[timeline span, draw=DefColor!76!black] (11.2,0.85) -- (13.2,0.85);

\node[timeline milestone] (m1) at (2.2,1.6) {基础概念};
\node[timeline milestone] (m2) at (4.5,1.6) {例题训练};
\node[timeline event] (e1) at (5.5,0) {};
\node[timeline milestone] (m3) at (7.5,1.6) {专题拓展};
\node[timeline milestone] (m4) at (10.0,1.6) {项目实践};
\node[timeline event, fill=LogicBg, draw=LogicColor!80!black] (e2) at (10.9,0) {};
\node[timeline milestone] (m5) at (12.2,1.6) {总结与发布};

\draw[timeline connector] (m1.south) -- (2.2,0.18);
\draw[timeline connector] (m2.south) -- (4.5,0.18);
\draw[timeline connector] (m3.south) -- (7.5,0.18);
\draw[timeline connector] (m4.south) -- (10.0,0.18);
\draw[timeline connector] (m5.south) -- (12.2,0.18);

\node[maki label, above=2pt of e1] {期中检查};
\node[maki label, above=2pt of e2] {验收};
\node[flow note, anchor=west] at (0.5,-1.45) {粗线可表示阶段, 圆点可表示硬节点, 矩形可表示主题段落.\\把文字改成章节、周次、版本号或历史年份都可以直接复用.};
\end{tikzpicture}
\captionof{figure}{课程安排或项目里程碑时间轴模板}
\end{center}

\clearpage

\section{概率与条件事件模板}

这一页适合概率论、统计推断或集合论导入部分. 左边保留原始样本空间, 右边展示“已知条件成立后”留下的有效区域, 这样既能讲事件之间的关系, 也能自然过渡到条件概率与归一化的概念.

\begin{center}
\begin{tikzpicture}[x=1cm,y=1cm]
\begin{scope}
  \draw[prob universe] (0,0) rectangle (5.8,4.4);
  \node[maki label, above left] at (0,4.4) {$\Omega$};
  \filldraw[prob event] (1.7,2.1) ellipse (1.15 and 1.55);
  \filldraw[prob event alt] (3.95,2.95) ellipse (1.0 and 0.78);
  \filldraw[prob event alt] (4.05,1.15) ellipse (1.0 and 0.78);
  \draw[prob condition] (2.95,0) -- (2.95,4.4);
  \foreach \x/\y in {1.1/3.25,1.45/2.55,1.85/1.85,1.4/1.15,2.0/2.9,3.7/3.0,4.25/2.65,3.7/1.1,4.3/0.95,4.7/1.35}{
    \node[prob outcome] at (\x,\y) {};
  }
  \node[prob label] at (1.7,3.85) {$A$};
  \node[prob label] at (4.05,3.65) {$B$};
  \node[prob label] at (4.08,1.9) {$C$};
  \node[prob label, anchor=west] at (3.1,4.08) {$A^c$};
  \node[flow note, anchor=west] at (0.15,-1.0) {左图保留全样本空间;\\虚线可表示被条件事件筛掉\\的区域边界.};
\end{scope}

\begin{scope}[xshift=8.2cm]
  \draw[prob universe] (0,0) rectangle (4.5,4.4);
  \filldraw[prob event] (1.55,2.2) ellipse (1.05 and 1.55);
  \foreach \x/\y in {1.1/3.05,1.45/2.45,1.7/1.85,1.35/1.1,1.95/2.85}{
    \node[prob outcome] at (\x,\y) {};
  }
  \node[prob label] at (1.55,3.85) {$A \mid A$};
  \node[maki annotation] at (2.75,2.2) {observe $A$\\discard $A^c$};
  \draw[flow arrow emph] (-1.55,2.2) -- (-0.25,2.2);
  \node[maki label] at (-0.9,2.65) {条件化};
  \node[flow note, anchor=west] at (0.0,-1.0) {右图可直接改成“给定 $B$”\\“给定测量值落在某区间内”\\或“过滤异常点后”的版本.};
\end{scope}
\end{tikzpicture}
\captionof{figure}{样本空间与条件事件更新模板}
\end{center}

\clearpage

\section{图论与网络结构模板}

这一页适合最短路、最小割、网络流、图搜索过程或结构关系图. 核心做法是把普通边、强调边、候选切边和搜索路径拆开, 这样同一张图既能承载“静态结构”, 也能承载“算法过程”.

\begin{center}
\begin{tikzpicture}[x=1cm,y=1cm]
\node[graph node accent] (s) at (0,1.9) {};
\node[graph node] (a) at (2.2,3.2) {};
\node[graph node] (b) at (4.6,2.6) {};
\node[graph node muted] (c) at (4.9,0.85) {};
\node[graph node] (d) at (2.0,0.1) {};
\node[graph node] (t) at (6.9,1.85) {};

\foreach \u/\v/\w in {
  s/a/2, s/d/3, a/b/2, a/d/1, d/b/3, d/c/2, b/c/2, b/t/1, c/t/3
}{
  \draw[graph edge] (\u) -- (\v) node[midway, graph label, fill=white] {\w};
}

\draw[graph edge strong] (s) -- (a);
\draw[graph edge strong] (a) -- (b);
\draw[graph edge strong] (b) -- (t);
\draw[graph edge cut] (d) -- (c);
\draw[graph walk] (s) to[bend right=14] (d);
\draw[graph walk] (d) to[bend left=8] (b);

\node[maki label, above left=2pt of s] {s};
\node[maki label, above=2pt of a] {a};
\node[maki label, above=2pt of b] {b};
\node[maki label, below=2pt of c] {c};
\node[maki label, below=2pt of d] {d};
\node[maki label, right=2pt of t] {t};

\node[flow note, anchor=west] at (8.0,2.95) {红色强调边: 当前主路径\\虚线边: cut / bottleneck candidate\\箭头: 搜索或松弛的访问轨迹};
\node[maki annotation] at (3.45,4.15) {shortest-path view};
\node[maki annotation] at (4.2,-1.0) {same graph, different overlays};
\end{tikzpicture}
\captionof{figure}{图搜索、最短路或网络流示意模板}
\end{center}

\end{document}
